\begin{abstract}
Los sistemas de potencia eléctrica (EPS) enfrentan crecientes demandas de eficiencia, flexibilidad e integración de energías renovables. Este artículo presenta una metodología integral para la optimización de EPS mediante convertidores DC-DC de alta eficiencia y técnicas de control digital avanzado. Se analizan topologías resonantes y de conmutación suave, junto con estrategias de control predictivo por modelo (MPC) y modulación de fase triple para minimizar pérdidas y mejorar la respuesta dinámica. Mediante modelado en estado-espacio, simulación en PSIM/MATLAB y validación experimental en prototipos de 1 kW, se logra eficiencia superior al 96% en modos bidireccionales. Los resultados demuestran reducción significativa en pérdidas de conmutación, estabilidad mejorada bajo cargas variables y optimización del flujo de potencia en microgrids DC. La integración de control digital permite implementación en tiempo real con DSP/FPGA, ofreciendo una solución reproducible y escalable para aplicaciones en vehículos eléctricos, almacenamiento energético y sistemas de distribución moderna. Se discuten limitaciones y direcciones futuras en control adaptativo y topologías de alta densidad de potencia.
\end{abstract}